\begin{scanpage}{0220}{202}
\[\left.\begin{aligned}
C_1&=-\frac1{2D}e^{2ikL}\cdot\ln\frac{i}{\gamma ka}\cdot\bigl[\psi_+^0\psi e^{ikL-ikz_2\cos\vartheta_0}-\psi_-^0 e^{-ikz_1\cos\vartheta_0}\bigr],\\
C_2&=-\frac1{2D}e^{2ikL}\cdot\ln\frac{i}{\gamma ka}\cdot\bigl[\psi_-^0\psi e^{ikL-ikz_1\cos\vartheta_0}-\psi_+^0 e^{-ikz_2\cos\vartheta_0}\bigr].
\end{aligned}\right\}\tag{32.14}\]
\sectiontitle{33}{A Vibrator Which is Short in Comparison with the Wavelength (a Passive Dipole)}
The theory of plane wave scattering by a thin cylindrical vibrator which is discussed in this chapter is based on a number of physical considerations. One good aspect of this theory is the fact that its precision increases with the length of the vibrator, since the current waves whose diffraction we are investigating are expressed more clearly, the longer the vibrator. However, one may also show that for short vibrators, the length of which is small in comparison with the wavelength, the equations derived by us have good precision.

It is clear that a vibrator which is short in comparison with the wavelength acts as a dipole, creating a fringing field
\[E_\vartheta=H_\varphi=-k^3p_z\frac{e^{ikR}}{kR}\sin\vartheta,\tag{33.01}\]
where the dipole moment $p_z$ may be calculated by solving the electrostatic problem. This dipole moment depends on the dimensions and shape of the vibrator. In accordance with [92], the dipole moment of a cylinder in a uniform electrostatic field $E_z$ equals
\[p_z=D(l)\left(\frac L2\right)^3\cdot E_z,\tag{33.02}\]
where $D(l)$ is a dimensionless function $l=L/2a$ which is shown in Figure 75 by the continuous curve. With $l\gg1$, one may calculate the function $D(l)$ by means of the asymptotic expansion
\[D(l)=\frac23\left(\frac1{\Omega_2}+\frac{0.977}{\Omega_2^3}+\cdots\right),\quad\Omega_2=2\left(\ln4l-\frac73\right).\tag{33.03}\]
\end{scanpage}
